% SPDX-License-Identifier: CC-BY-SA-4.0
% Author: Matthieu Perrin
% Part: <Nom de la partie>
% Section: <Nom de la section>
% Sub-section: <Nom de la sous-section>  % (facultatif, laisser vide si non utilisé)
% Frame: <Titre de la slide>

\begingroup

\begin{frame}{Conséquence : No Silver Bullet}

  \onAlertBlock[top=-3mm]{Conséquence pour l'analyse statique de programmes}{
    Il n'existe pas d'analyseur statique \alert{parfait}. Tout outil de vérification est soit
    \begin{description}
    \item[Incorrect :] il signale des "faux positifs", ou
    \item[Incomplet :] il rate certains bugs.
    \end{description}
  }

  \onExampleBlock[bottom=8mm]{Exemples de propriétés indécidables}{
    \begin{itemize}
    \item $P$ s'arrête sur une certaine entrée
    \item $P$ s'arrête sur toutes ses entrées
    \item $P$ ne déréférence pas le pointeur nul
    \item $P$ ne lève pas d'exception
    \item $P$ est un virus
    \item $P$ retourne le même résultat qu'un autre programme
    \item $P$ retourne un résultat correct par rapport à sa spécification
    \end{itemize}
  }  

  \footnoteref{Brooks. \emph{No Silver Bullet -- Essence and Accident in Software Engineering.} IFIP (1986)}
  \onImage[y=-2mm, x=40mm]{%
    width=2cm,
    title={Fred Brooks\footnote{Prix Turing 1999 pour ses contributions en architecture des ordinateurs, systèmes d'exploitation et logiciels}},
    licenselogo={\ccbysa},
    license={{\href{https://creativecommons.org/licenses/by-sa/3.0/}{CC BY-SA-3.0}} -- SD\&M, 2007 (\href{https://commons.wikimedia.org/wiki/File:Fred_Brooks.jpg}{Wikimedia})},
    img={Brooks.jpg}
  }

\end{frame}

\endgroup
